\documentclass[10pt,a4paper]{amsart}
\usepackage[margin=19mm]{geometry}
\usepackage{amsmath,amssymb,mathtools}
\usepackage[scr=rsfs,cal=euler]{mathalfa}
\usepackage{xfrac}
\usepackage[unicode,hidelinks,bookmarks=false]{hyperref}
\newcommand{\ie}{i.e.\ }
\newcommand{\cf}{cf.\ }
\newcommand{\resp}{respectively\ }
\newcommand{\Subsection}{\subsection}
\providecommand{\nobreakdash}{\nobreak\hbox{-}}
\setlength{\emergencystretch}{2em}
\multlinegap0pt
\usepackage{amssymb}
\usepackage{float}
\newtheorem{lemma}{Lemma}
\title{On Boolean Functions Maximizing Divergence for Correlated Sources}
\author{Jun Chen, Shun Watanabe,, Lei Yu}
\date{Source: arXiv:2607.28162v1 (CC BY 4.0)}
\begin{document}
\maketitle

\noindent\emph{Source and licence.} On Boolean Functions Maximizing Divergence for Correlated Sources. Version arXiv:2607.28162v1 (CC BY 4.0), distributed by arXiv under the Creative Commons Attribution 4.0 International licence, http://creativecommons.org/licenses/by/4.0/, as recorded in the arXiv licence metadata for this exact version and shown on the abstract page https://arxiv.org/abs/2607.28162v1. Full licence notice: \texttt{licenses/arxiv-2607.28162-CC-BY-4.0.txt}.\par\smallskip

\noindent\textit{Editorial scope.} Counted unit: the article's lemma lemma:inner-product-convexity (source0.tex lines 444--454). The article's complete original proof of it is delivered with it (source0.tex lines 456--513, one \texttt{proof} environment of 58 lines). No mathematics is written, repaired or re-derived: the statement and the proof are the article's own lines, in the article's own notation, and no retained line is substituted or re-flowed.  No further source-local context is needed: every cross-reference of every retained line resolves inside the two delivered spans.  Nothing was omitted from any retained span, so the package carries no omission record.  No retained line contains a citation, so the wrapper carries no bibliography and no bibliography entry.  No retained line contains a figure environment, an image-inclusion command or drawing code, so no image is delivered and the package declares no asset dependency.  Editorial formatting only, in addition: an AMS \texttt{amsart} preamble that restates the article's own declarations and macros used by the retained lines, namely the environments \texttt{lemma} on separate counters, together with the standard packages those lines need.  The retained environments print in this document as Lemma 1 in order of appearance, whereas the article prints the same items with the numbering of its own full document.  The complete native source of the article as distributed in the arXiv e-print for this exact version is bundled unchanged as \texttt{sources/206393/source0.tex}.
\begin{lemma} \label{lemma:inner-product-convexity}
Let $a = (a_1,\ldots,a_m), b= (b_1,\ldots,b_m) \in \mathbb{R}^m$ be unit vectors, i.e., $\sum_{i=1}^m a_i^2 =  \sum_{i=1}^m b_i^2  = 1$.
Then, for $\mu = (\mu_1,\ldots,\mu_m), \nu = (\nu_1,\ldots,\nu_m) \in (-1,1)^m$, we have\footnote{Note that, by the Cauchy-Schwarz inequality,
$\sum_{i=1}^m a_i b_i \mu_i \in (-1,1)$ and $\sum_{i=1}^m a_i b_i \nu_i \in (-1,1)$.}
\begin{align}
d\bigg( \frac{1 + \sum_{i=1}^m a_i b_i \mu_i}{2} \bigg\| \frac{1+ \sum_{i=1}^m a_i b_i \nu_i}{2} \bigg) 
\le \min\bigg[
\sum_{i=1}^m a_i^2 d\bigg( \frac{1+\mu_i}{2} \bigg\| \frac{1+\nu_i}{2} \bigg),~\sum_{i=1}^m b_i^2 d\bigg( \frac{1+\mu_i}{2} \bigg\| \frac{1+\nu_i}{2} \bigg)
\bigg].
\end{align}
\end{lemma}

\begin{proof}
For $t \in [0,1]$, let 
\begin{align*}
\psi(t) &:= d\bigg( \frac{1+ \sum_{i=1}^m a_i b_i(\nu_i + t(\mu_i-\nu_i))}{2} \bigg\| \frac{1+\sum_{i=1}^m a_i b_i \nu_i}{2}\bigg), \\
\phi(t) &:= \sum_{i=1}^m a_i^2 d\bigg( \frac{1+(\nu_i + t(\mu_i-\nu_i))}{2} \bigg\| \frac{1+\nu_i}{2}\bigg), \\
\varphi(t) &:= \sum_{i=1}^m b_i^2 d\bigg( \frac{1+(\nu_i + t(\mu_i-\nu_i))}{2} \bigg\| \frac{1+\nu_i}{2} \bigg).
\end{align*}
Then, the claim of the lemma is equivalently stated as 
\begin{align*}
\psi(1) \le \min[ \phi(1), \varphi(1)].
\end{align*}
Clearly, we have
\begin{align*}
\psi(0) = \phi(0) = \varphi(0) = 0.
\end{align*}
Moreover, we have
\begin{align*}
\psi^\prime(t) &= \frac{\sum_{i=1}^m a_i b_i (\mu_i-\nu_i)}{2} \ln \frac{\big(1+\sum_{i=1}^m a_i b_i (\nu_i + t(\mu_i-\nu_i))\big)\big(1-\sum_{i=1}^m a_i b_i \nu_i\big)}{\big(1-\sum_{i=1}^m a_i b_i (\nu_i + t(\mu_i-\nu_i))\big) \big(1+ \sum_{i=1}^m a_i b_i \nu_i\big)}, \\
\phi^\prime(t) &= \sum_{i=1}^m \frac{a_i^2(\mu_i-\nu_i)}{2} \ln \frac{\big( 1+ (\nu_i + t(\mu_i - \nu_i))\big)\big(1-\nu_i\big)}{\big(1- (\nu_i + t(\mu_i-\nu_i))\big)\big(1+\nu_i \big)}, \\
\varphi^\prime(t) &= \sum_{i=1}^m \frac{b_i^2(\mu_i-\nu_i)}{2} \ln \frac{\big(1+(\nu_i+t(\mu_i-\nu_i))\big)\big(1-\nu_i\big)}{\big(1-(\nu_i+t(\mu_i-\nu_i))\big)\big(1+\nu_i\big)}.
\end{align*}
In particular, we have
\begin{align*}
\psi^\prime(0) = \phi^\prime(0) = \varphi^\prime(0) = 0.
\end{align*}
Therefore, it suffices to show 
\begin{align} \label{eq:proof-inner-product-convexity-0}
\psi^{\prime\prime}(t) \le \min[ \phi^{\prime\prime}(t), \varphi^{\prime\prime}(t)]
\end{align}
for $t \in [0,1]$. Note that
\begin{align*}
\psi^{\prime\prime}(t) &= \frac{\big(\sum_{i=1}^m a_i b_i (\mu_i-\nu_i)\big)^2}{1 - \big(\sum_{i=1}^m a_i b_i (\nu_i + t(\mu_i - \nu_i))\big)^2}, \\
\phi^{\prime\prime}(t) &= \sum_{i=1}^m \frac{a_i^2 (\mu_i-\nu_i)^2}{1 - (\nu_i +t(\mu_i-\nu_i))^2}, \\
\varphi^{\prime\prime}(t) &= \sum_{i=1}^m \frac{b_i^2(\mu_i-\nu_i)^2}{1 - (\nu_i + t(\mu_i-\nu_i))^2}.
\end{align*}

Let $x_i := \frac{a_i(\mu_i-\nu_i)}{\sqrt{1-(\nu_i+t(\mu_i-\nu_i))^2}}$ and $y_i:= b_i \sqrt{1-(\nu_i + t(\mu_i-\nu_i))^2}$ for $i=1,\ldots,m$.
By the Cauchy-Schwarz inequality, we have
\begin{align*}
\bigg( \sum_{i=1}^m a_i b_i (\mu_i-\nu_i) \bigg)^2 &= \bigg( x_i y_i\bigg)^2 \\
&\le \bigg( \sum_{i=1}^m x_i^2 \bigg)\bigg( \sum_{i=1}^m y_i^2\bigg) \\
&= \bigg( \sum_{i=1}^m \frac{a_i^2 (\mu_i-\nu_i)^2}{1- (\nu_i +t(\mu_i-\nu_i))^2} \bigg) 
 \bigg( \sum_{i=1}^m b_i^2 \big(1- (\nu_i + t(\mu_i-\nu_i))^2\big) \bigg),
\end{align*}
which implies 
\begin{align}
\frac{\big(\sum_{i=1}^m a_i b_i (\mu_i - \nu_i)\big)^2}{\sum_{i=1}^m b_i^2\big(1 - (\nu_i + t(\mu_i-\nu_i))^2 \big)} 
 \le \sum_{i=1}^m \frac{a_i^2 (\mu_i - \nu_i)^2}{1 - (\nu_i + t(\mu_i-\nu_i))^2}. \label{eq:proof-inner-product-convexity-1}
\end{align}
Again by the Cauchy-Schwarz inequality, we have
\begin{align}
\sum_{i=1}^m b_i^2 \big( 1- (\nu_i + t(\mu_i-\nu_i))^2 \big) &= 1 - \sum_{i=1}^m b_i^2 (\nu_i + t(\mu_i-\nu_i))^2 \nonumber \\
&\le 1 - \bigg( \sum_{i=1}^m a_i b_i (\nu_i + t(\mu_i-\nu_i)) \bigg)^2. \label{eq:proof-inner-product-convexity-2}
\end{align}
Substituting \eqref{eq:proof-inner-product-convexity-2} into \eqref{eq:proof-inner-product-convexity-1} shows 
$\psi^{\prime\prime}(t) \le \phi^{\prime\prime}(t)$. By symmetry, we can also show $\psi^{\prime\prime}(t) \le \varphi^{\prime\prime}(t)$.
This completes the proof of \eqref{eq:proof-inner-product-convexity-0}.
\end{proof}

\end{document}
